\documentclass[11pt]{article}
\usepackage[a4paper,margin=1in]{geometry}
\usepackage{amsmath,amssymb,mathtools,bm}
\usepackage{enumitem,booktabs,array}
\usepackage{tcolorbox}
\usepackage{hyperref}
\usepackage[protrusion=true,expansion=false]{microtype}
\hypersetup{colorlinks=true,linkcolor=blue,urlcolor=blue,citecolor=blue}
\setlist[itemize]{topsep=2pt,itemsep=2pt}
\setlist[enumerate]{topsep=2pt,itemsep=2pt}
\newtcolorbox{keybox}{colback=blue!3!white,colframe=blue!40!black,boxrule=0.4pt,arc=1mm}
\newtcolorbox{alertbox}{colback=red!3!white,colframe=red!40!black,boxrule=0.4pt,arc=1mm}
\newtcolorbox{answerbox}{colback=green!3!white,colframe=green!40!black,boxrule=0.4pt,arc=1mm}
\newtcolorbox{drillbox}{colback=yellow!5!white,colframe=orange!60!black,boxrule=0.4pt,arc=1mm}
\newcommand{\EV}{\mathbb{E}}
\newcommand{\Var}{\mathrm{Var}}
\newcommand{\Cov}{\mathrm{Cov}}
\newcommand{\blank}[1]{\underline{\hspace{#1}}}

\title{E01-07: Friedman (1953, Essays in Positive Economics) \\
\large ``The Methodology of Positive Economics'' [BACKGROUND] --- Active Recall Workbook}
\author{Empirical Corporate Finance II --- Exam Prep}
\date{\today}
\begin{document}\maketitle

\begin{keybox}
\textbf{Paper in one sentence.} A positive economic theory should be judged only by the accuracy of its predictions --- not by the realism of its assumptions. Unrealistic-looking assumptions (perfect competition, profit maximization, rational expectations) are acceptable \emph{as-if} fictions so long as the theory predicts well.

\textbf{Placement.} Methodological foundation for neoclassical economics and modern corporate finance. Formalizes the ``as-if'' defense of rationality invoked by Alchian (1950). Central reference for how to evaluate empirical models.
\end{keybox}

\section*{Round 1: Complete Walk-Through}

\subsection*{1.1 The question}
What makes a good economic theory? Should the assumptions describe reality? Or should only predictions matter?

\subsection*{1.2 Friedman's position}
A theory is valuable if and only if it predicts observed phenomena accurately. Assumptions are \emph{instrumental}: they help generate predictions but do not need to be literally true.

\subsection*{1.3 The as-if defense}
\begin{itemize}
\item Firms may not consciously maximize profits, but markets select as if they do (Alchian 1950).
\item A billiard player may not solve optimization equations, but plays as if he did (Friedman's classic example).
\item The usefulness of the rational-actor model comes from its predictive accuracy, not cognitive realism.
\end{itemize}

\subsection*{1.4 Implications for economic practice}
\begin{itemize}
\item Models can abstract from realism (frictionless markets, perfect information) to sharpen prediction.
\item Empirical tests focus on implications, not assumptions.
\item Debates over whether assumptions are ``realistic'' are methodologically misplaced (Friedman's view).
\end{itemize}

\subsection*{1.5 Influence and critique}
\begin{itemize}
\item Foundational for modern empirical economics --- model is a tool, not a portrait.
\item Critics (Samuelson's revealed-preference; Musgrave 1981) argue that assumptions \emph{do} need to be consistent with domain.
\item Behavioural economics complicates Friedman: systematic bounds on rationality may break the as-if defense.
\item Modern practice: Friedman-style predictive tests coexist with structural models that do require realistic primitives.
\end{itemize}

\subsection*{1.6 Corporate-finance relevance}
\begin{itemize}
\item Justifies rational-agent models of capital structure, investment, and asset pricing, even when managers are clearly boundedly rational.
\item Reduced-form predictions-first empirical style dominates corporate-finance journals.
\item Behavioral corporate finance (Baker-Wurgler; Malmendier) challenges the as-if in settings where cognitive biases have large predictive content.
\end{itemize}

\subsection*{1.7 Caveats}
\begin{alertbox}
\begin{itemize}
\item The position is methodological, not formal; not universally accepted (Koopmans 1957, McCloskey 1985).
\item If predictions fail, assumptions come back into play for diagnosing the model.
\item In applied corporate finance, structural estimation requires taking assumptions seriously.
\end{itemize}
\end{alertbox}

\section*{Round 2: Level 1 Fill-in}
\begin{enumerate}
\item A theory is evaluated by the accuracy of its \blank{2cm}.
\item Assumptions are \blank{2cm}, not literal descriptions.
\item The ``\blank{2cm}'' defense extends Alchian's selection argument.
\item Example: a \blank{2cm} player plays as-if solving physics.
\item Behavioural economics challenges the as-if when biases are \blank{2cm}.
\end{enumerate}

\section*{Round 3: Level 2 Prompts}
\begin{enumerate}
\item State Friedman's methodological position and the as-if defense.
\item Relate Friedman (1953) to Alchian (1950).
\item Discuss the critique from Samuelson / Musgrave / behavioural economics.
\item What are the implications for reduced-form vs.\ structural empirical work?
\item In a corporate-finance empirical paper, when does the as-if defense seem strong vs.\ weak?
\end{enumerate}

\section*{Round 4: Minimal Scaffolding}
From a blank page: (i) question, (ii) predictive-accuracy position, (iii) as-if defense, (iv) implications, (v) critique.

\section*{Round 5: Blank Page}
Essay: the methodology of positive economics and its implications for empirical corporate finance.

\vspace{6pt}
\begin{drillbox}
\textbf{Speed Drill.} (1) Evaluation criterion. (2) Role of assumptions. (3) As-if example. (4) Relation to Alchian. (5) Main critique.
\end{drillbox}

\section*{Quick Reference \& Answer Key}
\begin{answerbox}
\begin{itemize}
\item \textbf{Criterion.} Predictive accuracy.
\item \textbf{Assumptions.} Instrumental, not literal.
\item \textbf{As-if.} Billiard player, profit max via selection.
\item \textbf{Critique.} Behavioural economics; assumptions matter when biases large.
\end{itemize}
\end{answerbox}

\begin{keybox}
\textbf{30-second exam pitch.} Friedman (1953) argues that a positive economic theory should be judged by the accuracy of its predictions, not by the realism of its assumptions. Assumptions are instrumental: a billiard player plays as if solving physics even though he does not; a firm behaves as if maximizing profit even if managers do not literally compute first-order conditions (Alchian 1950). This ``as-if'' defense justifies neoclassical and rational-agent models in applied economics and corporate finance, supporting reduced-form prediction-first empirical styles that dominate the field. Critics (Samuelson, Musgrave, and modern behavioural economics) argue that when assumptions are systematically violated in a domain, the as-if defense fails and predictions deteriorate. The paper is the background for why corporate finance routinely uses rational-agent tools while acknowledging (via Baker-Wurgler, Malmendier, etc.) behavioural departures. Methodologically, it frames how to read and critique empirical finance papers.
\end{keybox}

\end{document}
